\capituloUDA{Conclusiones y recomendaciones}{Cierre de resultados y criterios de aplicación}{Conclusiones\par Recomendaciones\par Limitaciones y trabajos futuros}

Este capítulo se completará al finalizar el análisis comparativo de escenarios. De acuerdo con las normas UDA, las conclusiones recogerán los aspectos más destacados del estudio y no se redactarán como viñetas.

\section{Conclusiones}

Las conclusiones deberán responder a la pregunta de investigación, contrastar la hipótesis y evaluar el cumplimiento de los objetivos. También deberán sintetizar los hallazgos sobre desempeño energético, confort térmico, estrategias pasivas, materiales sostenibles y pertinencia del modelo propuesto para el contexto insular.

\section{Recomendaciones}

Las recomendaciones deberán orientar la aplicación práctica de las estrategias identificadas, priorizando aquellas con mayor impacto técnico y factibilidad. También podrán incluir recomendaciones normativas, metodológicas y proyectuales para futuras edificaciones públicas en Galápagos.

\section{Limitaciones y trabajos futuros}

Las principales limitaciones previstas se relacionan con el acceso a información técnica actualizada, posibles cambios del diseño durante la ejecución, disponibilidad de datos climáticos y precisión de los supuestos de modelado. Los trabajos futuros podrán ampliar la validación con mediciones en obra, monitoreo post-ocupación y comparación con otros edificios públicos del archipiélago.
